\chapter{方差分析和回归分析}

\section{方差分析}

方差分析(Analysis of Variance, AoV or ANOVA)由著名英国统计学家R. A. Fisher在1923年提出，又称F检验，用于两个及两个以上样本均数差别的显著性检验。

\begin{definition}{}{}
	试验指标：要考察的指标；

	因素：影响试验指标的条件；

	因素的水平：因素所处的状态；
	
	单因素试验：在一项试验中只有一个因素在改变，

	多因素试验：多于一个因素在改变。
\end{definition}

基本假设：
\begin{compactenum}
	\item 各观测值相互独立，每组观测值来自正态总体
	\item 各总体方差相等
\end{compactenum}

\paragraph{变差分解}

组数$s$，总例数$N = \textstyle\sum_{j=1}^sn_j$

\begin{definition}{总变差}{total sum of square}
	总变差的大小可用总偏差平方和表示，这反映了所有观测值之间总的变差程度。
	\begin{subequations}
		\begin{align}
			S_{\mathrm T} := {}&\sum_{j=1}^s\sum_{i=1}^{n_j}(X_{ij} - \avg X)^2\\\notag
			= {}&\sum_{i,j}X_{ij}^2 - \frac1N\biggp{\sum_{i,j}X_{ij}}^2,
		\end{align}
	\end{subequations}
	总自由度$\nu_{\mathrm T} = N - 1$，总均方$\mathrm{MS}_{\mathrm T} = S_{\mathrm T}/\nu_{\mathrm T}.$
\end{definition}

变异程度除与偏差平方和$S$有关外，还与其自由度$\nu$有关，由于各部分自由度不相等，因此各部分偏差平方和不能直接比较，须将各部分偏差平方和除以相应自由度，其比值称为均方差 ，简称均方(mean square, MS)。

\begin{definition}{组间变差}{error sum of square}
	各处理组的样本均值大小不等, 这种变差称为组间变差%，其大小可用各组均值$X_{\cdot j}$与总均值$X$的离均差平方和表示。记作$S_{\mathrm A}$
	\begin{subequations}
		\begin{align}
			S_{\mathrm A} := {}&\sum_{j=1}^sn_j(\avg X_{\cdot j} - \avg X)^2\\\notag
			= {}&\sum_{j=1}^s\frac1{n_j}\biggp{\sum_{i=1}^{n_j}X_{ij}}^2 - \frac1N\biggp{\sum_{i,j}X_{ij}}^2,
		\end{align}
	\end{subequations}
	组间自由度$\nu_{\mathrm A} = s - 1$，组间均方$\mathrm{MS}_{\mathrm A} = S_{\mathrm A}/\nu_{\mathrm A}.$
\end{definition}
组间变差存在的原因
\begin{compactitem}
	\item 随机误差，包括个体变差和测量误差
	\item 处理因素的不同水平可能对试验结果有影响。
\end{compactitem}
\begin{definition}{组内变异}{treatment sum of square}
	在同一处理组内，虽然各受试对象接受的处理相同，但测量值之间仍不同，这种变异称为组内变异(误差)
	\begin{align}
		S_{\mathrm E} := {}&\sum_{i,j}(X_{ij} - \avg X_{\cdot j})^2%\\\notag
		%={}&\sum_{j=1}^s(n_j-1)S_i^2.
	\end{align}
	反映了随机误差的影响。

	组内自由度$\nu_{\mathrm E} = N - s$，组内均方$\mathrm{MS}_{\mathrm E} = S_{\mathrm E}/\nu_{\mathrm E}.$
\end{definition}
偏差平方和与自由度具有可加性
\begin{subequations}
	\label{eq:ps-anova-decomposition}
	\begin{align}
		S_{\mathrm T} &= S_{\mathrm A} + S_{\mathrm E},\\
		\nu_{\mathrm T} &= \nu_{\mathrm A} + \nu_{\mathrm E}.
	\end{align}
\end{subequations}
单因素方差分析检验各组的总体均值是否相等：
\[
	H_0:\mu_1 = \cdots = \mu_s,\qquad
	H_1:\text{至少有两个总体均值不相等}.
\]
检验统计量为组间均方与组内均方之比。在 $H_0$ 成立时，
\begin{equation}
	F := \frac{\mathrm{MS}_{\mathrm A}}{\mathrm{MS}_{\mathrm E}}
	= \frac{S_{\mathrm A}/(s - 1)}{S_{\mathrm E}/(N - s)}
	\sim \FDist(s - 1,N - s).
\end{equation}
记 $F_\alpha(s - 1,N - s)$ 为该分布的上侧 $\alpha$ 分位数。在显著性水平 $\alpha$ 下，若观测到的 $F > F_\alpha(s - 1,N - s)$，则拒绝 $H_0$。
\paragraph{双因素检验}
设因素 $\mathrm A$ 有 $r$ 个水平 $\mathrm A_1,\ldots,\mathrm A_r$，因素 $\mathrm B$ 有 $s$ 个水平 $\mathrm B_1,\ldots,\mathrm B_s$。对每个水平组合 $(\mathrm A_i,\mathrm B_j)$ 都做 $t$ 次试验，称为等重复试验。

设$X_{ijk} \sim \NormalDist(\mu_{ij},\sigma^2)$，$i, j, k $分别为因素 $\mathrm A,\mathrm B$ 的水平编号和重复试验的编号。
则
总平均
\[
	\mu := \frac1{rs}\sum_{i,j}\mu_{ij},
\]
水平 $\mathrm A_i,\mathrm B_j$ 的效应分别为
\begin{equation}\label{eq:ps-main-effects}
	\alpha_i = \mu_{i\cdot} - \mu,\quad\beta_j = \mu_{\cdot j} - \mu,
\end{equation}
交互效应
\begin{equation}\label{eq:ps-interaction-effect}
	\gamma_{ij} = \mu_{ij} - \mu_{\cdot j} - \mu_{i\cdot} + \mu.
\end{equation}
可见
\[
	\sum_{i=1}^r\alpha_i = \sum_{j=1}^s\beta_j = 0,
	\qquad \sum_{i=1}^r\gamma_{ij} = 0\ (\forall j),
	\qquad \sum_{j=1}^s\gamma_{ij} = 0\ (\forall i).
\]
\begin{compactitem}
	\item $S_{\mathrm T}$为总变差，自由度$rst - 1$
	\item $S_{\mathrm E}$误差平方和，自由度$rs(t - 1)$
	\item $S_{\mathrm A},S_{\mathrm B}$为因素 $\mathrm A,\mathrm B$ 的效应平方和，自由度为$ r - 1, s - 1$
	\item $S_{\mathrm A\times\mathrm B}$为 $\mathrm A,\mathrm B$ 的交互效应平方和，自由度为 $(r - 1)(s - 1)$
\end{compactitem}
\begin{align}
	\notag
	S_{\mathrm T} = {}& \sum_{i=1}^r \sum_{j=1}^s \sum_{k=1}^t(X_{i j k} - \avg{X})^2 \\
	\notag
	= {}& \sum_{i=1}^r \sum_{j=1}^s \sum_{k=1}^t[(X_{i j k} - \avg{X}_{i j \cdot}) + (\avg{X}_{i\cdot\cdot} - \avg{X}) + (\avg{X}_{\cdot j \cdot} - \avg{X}) +\\
	\notag
	&\quad(\avg{X}_{i j \cdot} - \avg{X}_{i \cdot\cdot} - \avg{X}_{\cdot j \cdot} + \avg{X})]^2 \\
	\notag
	= {}& \sum_{i=1}^r \sum_{j=1}^s \sum_{k=1}^t(X_{i j k} - \avg{X}_{i j \cdot})^2 + s t \sum_{i=1}^r(\avg{X}_{i\cdot\cdot} - \avg{X})^2 + \\
	\notag
	& r t \sum_{j=1}^s(\avg{X}_{\cdot j \cdot} - \avg{X})^2 + t \sum_{i=1}^r \sum_{j=1}^s(\avg{X}_{i j \cdot} - \avg{X}_{i\cdot\cdot} - \avg{X}_{\cdot j \cdot} + \avg{X})^2\\
	= {}&S_{\mathrm E} + S_{\mathrm A} + S_{\mathrm B} + S_{\mathrm A\times\mathrm B}.
\end{align}
例如，对因素 $\mathrm A$ 做假设检验
\[
	H_0:\alpha_1 = \cdots = \alpha_r = 0,\quad H_1:\text{至少一个 }\alpha_i \ne 0
\]
在 $H_0$ 成立且每组重复数 $t > 1$ 时，检验统计量服从
\begin{equation}\label{eq:ps-two-way-f-test}
	F := \biggdivision{\frac{S_{\mathrm A}}{r - 1}}{\frac{S_{\mathrm E}}{rs(t - 1)}} \sim \FDist\bigp{r - 1,rs(t - 1)}.
\end{equation}

\section{回归分析}
%如果变量$x$和$y$从相关图中可以看出它们之间大致形成一种直线关系，我们就可在相关图上求出一条与各点最相配合的直线。一元线性回归方程$y=a+bx$

%设$(x_1,y_1),\ldots,(x_n,y_n) $是$ (x, y) $的一组观测值，对每个样本观测值$(x_i,y_i)$考虑与其回归值$a+bx_i$的离差为$y_i-a-bx_i$，离差平方和
设$X,Y$的二阶矩有限且方差均非零。用$g(X)$估计$Y$，离差$Y - g(X)$，均方误差就是离差平方的期望
\[
	Q := \E\bigb{(Y - g(X))^2}.
\]
最小二乘法所求的是均方误差意义下的最优线性预测
\begin{equation}\label{eq:ps-least-squares-predictor}
	(\alpha,\beta) = \argmin_{a,b}\E\bigb{(Y - (a + bX))^2}
\end{equation}
实际上就是求二元函数
\begin{align*}
	f(a,b) :=&{}\E\bigb{(Y - (a + bX))^2}\\
	=&{}\E(X^2)b^2 + 2\E(X)ab + a^2 - 2\E(XY)b - 2\E(Y)a + \E(Y^2).
\end{align*}
的最小值点，则在最小值点处有
\[
	\pv fa = 0,\enspace\pv fb = 0.
\]
解得
\begin{equation}
	\begin{cases}
		\beta = \frac{\E(XY) - \E(X)\E(Y)}{\E(X^2) - \E^2(X)} \equiv \frac{\Cov(X,Y)}{\Var(X)},\\
		\alpha = \E(Y) - \beta\E(X)
	\end{cases}
\end{equation}
误差的期望和方差
\begin{align*}
	\E[Y - (\alpha + \beta X)] &= \E(Y) - \alpha - \beta\E(X) = 0;\\
	\Var[Y - (\alpha + \beta X)] &= \Var(Y - \beta X)\\
	&= \Var(Y) + \beta^2\Var(X) - 2\beta\Cov(Y,X)\\
	&=%\Var(Y)-\frac{\Cov^2(X,Y)}{\Var(X)}=
	\bigb{1 - \Corr^2(X,Y)}\Var(Y).
\end{align*}
\paragraph{相关性检验}
对任意两个变量的一组观察值$ {(x_i, y_i)} $都可以用最小二乘法形式上求得$ y$对$ x $的回归方程，如果$ y $与$ x $没有线性相关关系, 这种形式的回归方程就没有意义。因此需要考察$ y $与$ x $间是否确有线性相关关系，这就是回归效果的检验问题。

固定$n > 2$个不全相同的自变量值，设各观测误差独立、同方差。正态一元线性回归模型为
\begin{equation}\label{eq:ps-simple-linear-regression}
	y = a + bx + \epsilon,\enspace\epsilon \sim \NormalDist(0,\sigma^2).
\end{equation}
因随机因素引起的误差称为残差平方和
\begin{equation}\label{eq:ps-residual-sum-squares}
	Q_{\mathrm e} = \sum_{i=1}^n(y_i - \hat y_i)^2
\end{equation}
则
\begin{equation}\label{eq:ps-regression-error-variance}
	\frac{Q_{\mathrm e}}{\sigma^2} \sim \ChiSqDist(n - 2) \implies \hat\sigma^2 = \frac{Q_{\mathrm e}}{n - 2},
\end{equation}
检验假设 
\[
	H_0:b = 0,\quad H_1:b \neq 0,
\]
$\hat b$满足 
\begin{equation}\label{eq:ps-slope-distribution}
    \begin{aligned}
		\hat b &\sim \NormalDist\biggp{b,\frac{\sigma^2}{L_{xx}}},\\
		L_{xx} &:= \sum_{i=1}^n(x_i - \avg x)^2.
    \end{aligned}
\end{equation}
如果原假设成立，
\begin{align*}
	\frac{\hat b\sqrt{L_{xx}}}{\sigma} &\sim \NormalDist(0,1),\\
	\frac{Q_{\mathrm e}}{\sigma^2} &\sim \ChiSqDist(n - 2).
\end{align*}
检验统计量 
\begin{equation}
	t := \hat b\sqrt{\frac{(n - 2)L_{xx}}{Q_{\mathrm e}}} \sim \TDist(n - 2).
\end{equation}
拒绝域$\abs t \geqslant t_{\alpha/2}(n - 2).$
\paragraph{线性化回归}
在许多实际问题中，两个变量之间并不一定是线性关系，而是某种曲线关系，应该用曲线来拟合。可以进行适当的变量代换，把它线性化，这样就把一个非线性回归问题化为线性回归问题而得以解决。
\paragraph{多线性回归}
随机变量$ Y $往往与多个变量$ x1, x2, \ldots , xp $有关。对自变量$ x1, x2, \ldots, xp $的一组确定的值，$Y $有它的分布
\begin{equation}\label{eq:ps-multiple-linear-regression}
	Y = b_0 + \sum_{i=1}^pb_ix_i + \epsilon,\enspace\epsilon \sim \NormalDist(0,\sigma^2).
\end{equation}
